% ============================================================
%  Resumen
% ============================================================
\begin{center}
    \textbf{\huge Resumen}
\end{center}
\vspace{1cm}

En la actualidad, la seguridad eléctrica de los equipos de uso doméstico y
comercial constituye un requisito fundamental para su comercialización. Sin
embargo, los métodos tradicionales de ensayo, basados en la intervención manual
del operario, pueden generar variaciones en los resultados, extender los tiempos
de prueba y aumentar los riesgos para el personal y los dispositivos bajo
evaluación.

Este trabajo presenta el desarrollo de un sistema integral, conformado por un
dispositivo de ensayos y una interfaz de software, destinado a verificar en final
de línea el cumplimiento de la norma IEC~60335-1 en aparatos electrodomésticos:
consumo, corriente de fuga activa y pasiva, resistencia de puesta a tierra y
rigidez dieléctrica. El dispositivo se basa en un microcontrolador ESP32 que lee
medidores PZEM-004T y en relés y contactores que conmutan los circuitos de ensayo;
la interfaz, desarrollada en Python y PyQt5, configura y secuencia los ensayos,
emite el veredicto, registra los resultados en una base de datos y los exporta.

Durante el desarrollo se alinearon los criterios de aprobación con los ensayos de
rutina del Anexo~A y las cláusulas~10 y~16 de la norma, se escribió el firmware
del microcontrolador y se integró y verificó el sistema en un banco de pruebas
real. El software se probó con hardware real en dos equipos: en un equipo
completo provisto por SYGlab y esquematizado por la autora, donde su versión
inicial ejecutó los ensayos de seguridad y registró 86 verificaciones con
patrones, y en un prototipo propio, donde la versión final se verificó de punta a
punta en el ensayo de consumo sobre cargas reales. Los cambios de la versión
final en los ensayos de fuga, puesta a tierra y rigidez dieléctrica se
verificaron con hardware simulado. Se documentan además las
dificultades encontradas ---entre ellas un problema de compatibilidad
electromagnética en el control de los relés--- y las soluciones adoptadas.

\newpage
